\subsection{Hiperparametrai}
\label{subsec:hiperparametrai}

Visiems eksperimentams bendrosios mokymo hiperparametrų reikšmės yra fiksuotos: paketo dydis (angl.~\textit{batch size}) – $64$, epochų skaičius – $20$, mokymosi greitis (angl.~\textit{learning rate}) – $10^{-3}$, paklaidos funkcija – \textit{Cross Entropy}, optimizavimo algoritmas – \textit{Adam}, branduolio dydis – $3$, telkimo dydis – $2$. Šios reikšmės keičiamos tik tuose tyrimuose, kurie skirti būtent tam hiperparametrui įvertinti. Bendrieji hiperparametrai surinkti \ref{tab:hiperparametrai_bendri} lentelėje, o tyrimuose varijuojamos reikšmės – \ref{tab:hiperparametrai_tyrimai} lentelėje.

\begin{table}[H]
    \centering
    \caption{Bendrosios (fiksuotos) hiperparametrų reikšmės abiem tinklams.}
    \begin{tabular}{|l|l|}
        \hline
        Hiperparametras & Reikšmė \\
        \hline
        Paketo dydis (\textit{batch size}) & $64$ \\ \hline
        Epochų skaičius & $20$ \\ \hline
        Mokymosi greitis (\textit{learning rate}) & $10^{-3}$ \\ \hline
        Optimizavimo algoritmas & \textit{Adam} \\ \hline
        Paklaidos funkcija & \textit{Cross Entropy} \\ \hline
        Aktyvacijos funkcija & \textit{ReLU} \\ \hline
        Branduolio dydis (\textit{kernel}) & $3$ \\ \hline
        Telkimo dydis (\textit{pool}) & $2$ \\ \hline
        \texttt{padding} konvoliucijose & $1$ \\ \hline
        \textit{Dropout} tikimybė (bazinė) & $0{,}0$ \\ \hline
        Paketinis normalizavimas (bazinė) & išjungta \\ \hline
        Pradinė \textit{seed} reikšmė & $42$ \\ \hline
    \end{tabular}
    \label{tab:hiperparametrai_bendri}
\end{table}

\begin{table}[H]
    \centering
    \caption{Tyrimuose varijuojami hiperparametrai ir jų reikšmės.}
    \begin{tabular}{|l|l|l|}
        \hline
        Tyrimas & Vaizdų variantai & Laiko eilučių variantai \\
        \hline
        Architektūra (filtrai) & $(16,\,32)$; $(32,\,64,\,128)$; $(64,\,128,\,256)$ & $(32)$; $(64,\,128)$; $(64,\,128,\,256)$ \\ \hline
        \textit{Dropout} tikimybė & \multicolumn{2}{c|}{$0{,}0$; $0{,}25$; $0{,}5$; $0{,}75$} \\ \hline
        Paketinis normalizavimas & \multicolumn{2}{c|}{išjungta; įjungta} \\ \hline
        Aktyvacijos funkcija & \multicolumn{2}{c|}{\textit{ReLU}, \textit{LeakyReLU}, \textit{Tanh}, \textit{ELU}} \\ \hline
        Optimizavimo algoritmas & \multicolumn{2}{c|}{\textit{Adam}, \textit{SGD}, \textit{RMSprop}, \textit{AdamW}} \\ \hline
    \end{tabular}
    \label{tab:hiperparametrai_tyrimai}
\end{table}


Atsitiktinių skaičių generatoriaus pradinė reikšmė (angl.~\textit{seed}) buvo fiksuota ($seed = 42$) duomenų dalijime ir poaibių formavime, kad eksperimentų rezultatai būtų atkartojami.
